% Generated draft: requires mathematical and source review.
\documentclass[11pt]{article}
\usepackage[margin=25mm]{geometry}
\usepackage{fontspec}
\setmainfont{Latin Modern Roman}
\usepackage{amsmath,amssymb,mathtools,unicode-math}
\setmathfont{Latin Modern Math}
\usepackage{enumitem,needspace,xurl,hyperref}
\setlength{\parindent}{0pt}
\setlength{\parskip}{5pt}
\setlength{\emergencystretch}{3em}
\newcommand{\N}{\mathbb{N}}
\newcommand{\Z}{\mathbb{Z}}
\newcommand{\Q}{\mathbb{Q}}
\newcommand{\R}{\mathbb{R}}
\newcommand{\C}{\mathbb{C}}
\newcommand{\F}{\mathbb{F}}
\newcommand{\A}{\mathbb{A}}
\newcommand{\PP}{\mathbb{P}}
\newcommand{\E}{\mathbb{E}}
\newcommand{\eps}{\varepsilon}
\newcommand{\dd}{\,\mathrm{d}}
\begin{document}
\section*{UnsolvedMath ID: 6700003}
\textbf{Source problem code:} catalogue entry
\subsection{Definitions and mathematical statement}
For a smooth Riemannian manifold $(M^n,g)$, scalar curvature is the trace $\operatorname{Sc}_g=\operatorname{tr}_g\operatorname{Ric}_g$. For a prescribed scalar threshold $\sigma$, the condition in the source is the strict pointwise inequality
\[
\operatorname{Sc}_g(x)>\sigma\qquad(x\in M).
\]
For a constant threshold one can denote the metrics satisfying it by $\mathcal R_{>\sigma}(M)$. The source's question, with spacing normalized, is: what are the geometries of individual manifolds with this curvature inequality?

Describe the geometric possibilities and restrictions for individual Riemannian manifolds satisfying this bound. The item does not state a comparison theorem, geometric invariant to estimate, manifold category, dimension range, completeness or compactness assumption, or a proposed classification. It is an umbrella program and is not silently replaced by a question about the topology of the space $\mathcal R_{>\sigma}(M)$.

The supplied summary points to comparison and rigidity results for scalar-curvature bounds, but supplies no full hypotheses. In particular, a spherical lower bound alone is not here asserted to force a round metric, and the summary's references to tori/aspherical manifolds do not establish an unrestricted theorem for every such manifold. These are related areas of partial progress, while the exact source condition remains $\operatorname{Sc}>\sigma$, not $\operatorname{Sc}\geq\sigma$.
\subsection{Short English statement}
Describe what geometry is possible for an individual manifold with scalar curvature strictly greater than a prescribed threshold. The source poses a broad research program without a specific comparison or classification assertion.
\subsection{Sources}
\begin{itemize}
\item Misha Gromov, \emph{101 Questions, Problems and Conjectures around Scalar Curvature} (2017), printed tag [?3], occurrence 3, PDF page 1, \url{https://www.ihes.fr/~gromov/wp-content/uploads/2018/08/101-problemsOct1-2017.pdf}; accessed 2026-07-29, PDF-text-normalized extraction. The provenance describes the author-hosted PDF as incomplete.
\item Related references named in the supplied summary: Llarull; Gromov--Lawson; Brendle--Marques--Neves; and Gromov, \emph{Four Lectures on Scalar Curvature} (2023). No precise theorem citations or full hypotheses are given for these comparisons.
\item Local registry \texttt{data/unsolvedmath/problems.json}, permanent UnsolvedMath ID $6700003$, source code catalogue entry, \emph{Scalar Curvature Question [?3]}. Dataset attribution: Ulam AI and the UnsolvedMath Contributors.
\end{itemize}
\end{document}
